\documentclass{article}
\title{Cryoprotective Agent Study}
\author{Research Team}
\date{2026}

\begin{document}
\maketitle

\section{Introduction}

This paper investigates cryoprotective agents and their mechanisms of action.
We focus on membrane permeability and thermal dynamics.

\subsection{Background}

Cryoprotection is the process of protecting biological material from damage caused by freezing.
Two classes exist: penetrating and non-penetrating agents.

\subsubsection{Penetrating Agents}

Penetrating agents such as DMSO and glycerol cross cell membranes.
They reduce the concentration of electrolytes during freezing.

\subsection{Research Questions}

This study addresses three primary questions:
(1) What concentration thresholds minimize cytotoxicity?
(2) How does cooling rate interact with agent concentration?

\section{Methods}

\subsection{Experimental Design}

Samples were prepared in triplicate at five concentration levels.
Controls used phosphate-buffered saline without cryoprotective agents.

\subsection{Protocol}

Cooling was performed at 1°C/min from 20°C to -80°C.

\section{Results}

Preliminary data shows 85\% viability at optimal DMSO concentrations.

\section{Discussion}

The synergistic effect of combined penetrating and non-penetrating agents
warrants further investigation.

\end{document}
